% !TEX root = ../main.tex
\chapter{The Simple Regression Model}
\label{ch:slr}

Econometrics begins with a deceptively simple question: how does one variable move with another, \emph{on average}, once we admit that the relationship is not exact? Suppose we believe that wages rise with education, that consumption rises with income, or that a crop yield rises with rainfall. In each case the relationship is real but noisy --- two people with the same education earn different wages. The simple regression model is the smallest formal object that captures this idea: a straight line for the average, plus an error term for everything else.

This chapter introduces that model, defines the ordinary least squares (OLS) estimator that fits it, and asks the two questions we will ask of \emph{every} estimator in this course: is it correct on average (unbiased), and how precise is it (variance)? The answers here, in the one-regressor case, set the template for the multiple regression model in Chapter~\ref{ch:mlr}.

\section{The Population Model}

\defn{Simple Linear Regression Model}{
The \emph{simple linear regression (SLR) model} relates a dependent variable $y$ to a single explanatory variable $x$ through
\[
y = \beta_0 + \beta_1 x + u,
\]
where $\beta_0$ is the intercept, $\beta_1$ the slope, and $u$ the error (or disturbance) term collecting all factors other than $x$ that affect $y$.
}

The vocabulary is worth fixing once, because different fields use different names for the same objects:
\begin{itemize}
    \item $y$: \emph{dependent} variable, explained variable, response, or regressand;
    \item $x$: \emph{independent} variable, explanatory variable, control, or regressor;
    \item $u$: \emph{error} term or disturbance.
\end{itemize}

The slope $\beta_1$ is the object of interest: it is the change in $y$ associated with a one-unit change in $x$, \emph{holding the error fixed}, i.e.\ $\Delta y = \beta_1 \, \Delta x$ when $\Delta u = 0$. For this to have a causal reading we need the error to be unrelated to $x$, which we now make precise.

Normalizing $\E{u}=0$ is free: any nonzero mean of $u$ can be absorbed into the intercept $\beta_0$. The substantive assumption is the next one.

\defn{Zero Conditional Mean}{
The error has \emph{zero conditional mean} if
\[
\E{u \given x} = 0 \qquad \text{for every value of } x.
\]
Equivalently, the average of the unobserved factors does not vary with $x$.
}

Under zero conditional mean, taking the expectation of the model conditional on $x$ gives the \emph{population regression function (PRF)},
\[
\E{y \given x} = \beta_0 + \beta_1 x,
\]
which says that the conditional mean of $y$ is linear in $x$. This is what OLS will try to recover from a sample.

\section{Ordinary Least Squares}

Zero conditional mean implies two \emph{population moment conditions}: $\E{u}=0$ and $\E{xu}=0$ (the latter follows because $\E{u\given x}=0 \Rightarrow \Cov{x}{u}=0$). OLS is the method that imposes the sample analogues of these two conditions and solves for the two unknowns $\beta_0,\beta_1$.

\defn{OLS Estimators}{
Given a sample $\{(x_i,y_i):i=1,\dots,n\}$, the OLS estimators are
\[
\hb_1 = \frac{\sumi (x_i-\bar x)(y_i-\bar y)}{\sumi (x_i-\bar x)^2}
      = \frac{\sumi (x_i-\bar x)\,y_i}{\sumi (x_i-\bar x)^2},
\qquad
\hb_0 = \bar y - \hb_1 \bar x.
\]
The fitted (sample regression) function is $\hat y_i = \hb_0 + \hb_1 x_i$, and the residual is $\hat u_i := y_i - \hat y_i$.
}

The two equalities for $\hb_1$ coincide because $\sumi (x_i-\bar x)\bar y = \bar y \sumi(x_i-\bar x) = 0$; more generally $\sumi c\,(x_i-\bar x)=0$ for any constant $c$, a trick we will reuse constantly.

The slope can also be written in terms of moments,
\[
\hb_1 = \frac{\Cov{x}{y}}{\Var{x}} = \widehat\rho_{xy}\cdot \frac{\widehat\sigma_y}{\widehat\sigma_x},
\]
which shows that the sign of the slope is the sign of the sample correlation between $x$ and $y$.

\thm{Algebraic Properties of OLS}{
The OLS fit satisfies three identities by construction:
\begin{enumerate}
    \item $\sumi \hat u_i = 0$ \quad (residuals sum to zero, hence have zero sample mean);
    \item $\sumi x_i \hat u_i = 0$ \quad (regressor and residuals have zero sample covariance);
    \item $\bar y = \hb_0 + \hb_1 \bar x$ \quad (the point of means lies on the regression line).
\end{enumerate}
}

These are not assumptions about the world; they hold in \emph{every} sample because they are exactly the first-order conditions of the least-squares problem $\min_{b_0,b_1}\sumi (y_i - b_0 - b_1 x_i)^2$.

\section{Goodness of Fit}

How much of the variation in $y$ does the line explain? Decompose the total variation into an explained part and a residual part.

\defn{Sums of Squares and $R^2$}{
\[
\ba
\text{SST} &= \sumi (y_i-\bar y)^2 &&\text{(total)}\\
\text{SSE} &= \sumi (\hat y_i-\bar y)^2 &&\text{(explained)}\\
\text{SSR} &= \sumi (y_i-\hat y_i)^2 &&\text{(residual)}
\ea
\]
With $\text{SST}=\text{SSE}+\text{SSR}$, the \emph{coefficient of determination} is
\[
R^2 = \frac{\text{SSE}}{\text{SST}} = 1 - \frac{\text{SSR}}{\text{SST}} \in [0,1].
\]
}

$R^2$ is the fraction of the sample variation in $y$ explained by the model. Two cautions worth internalizing now, because they recur in the multiple regression model (Chapter~\ref{ch:mlr}):
\begin{itemize}
    \item $R^2$ equals the squared sample correlation between $y_i$ and its fitted value $\hat y_i$; it is a measure of fit, \emph{not} of causality.
    \item A low $R^2$ does not invalidate the regression. We may estimate the slope $\beta_1$ very precisely even when most of the variation in $y$ is left to the error.
\end{itemize}

\section{Units of Measurement and Functional Form}

The linear model is more flexible than it looks: by taking logs of $y$, of $x$, or of both, the \emph{same} OLS machinery delivers different and often more natural interpretations of the slope. The four standard cases are summarized below; $\%\Delta$ denotes a percentage change.

\begin{center}
\renewcommand{\arraystretch}{1.6}
\begin{tabular}{@{}llll@{}}
\toprule
Model & Dep.\ var. & Indep.\ var. & Interpretation of $\beta_1$ \\
\midrule
Level--level & $y$       & $x$       & $\Delta y = \beta_1\,\Delta x$ \\
Level--log   & $y$       & $\log x$  & $\Delta y = (\beta_1/100)\,(\%\Delta x)$ \\
Log--level   & $\log y$  & $x$       & $\%\Delta y \approx (100\,\beta_1)\,\Delta x$ \\
Log--log     & $\log y$  & $\log x$  & $\%\Delta y = \beta_1\,(\%\Delta x)$ \quad (elasticity) \\
\bottomrule
\end{tabular}
\end{center}

The log--log slope is an elasticity, and logged slopes are invariant to rescaling the units of the logged variable --- two reasons logs are ubiquitous in applied work. We return to functional form systematically in Chapter~\ref{ch:further}.

\section{Expected Value and Variance of the OLS Estimators}

We now ask the two questions that organize the rest of the course. To answer them we need a list of assumptions; the first four deliver unbiasedness, and the fifth (homoskedasticity) delivers the simple variance formula.

\asm{SLR.1--SLR.5 (Gauss--Markov Assumptions, Simple Regression)}{
\begin{description}
    \item[SLR.1 — Linear in parameters.] The population model is $y = \beta_0 + \beta_1 x + u$.
    \item[SLR.2 — Random sampling.] $\{(x_i,y_i):i=1,\dots,n\}$ is a random sample from the population, so $y_i = \beta_0 + \beta_1 x_i + u_i$.
    \item[SLR.3 — Sample variation in $x$.] The $x_i$ are not all equal: $\text{SST}_x := \sumi (x_i-\bar x)^2 > 0$.
    \item[SLR.4 — Zero conditional mean.] $\E{u \given x} = 0$.
    \item[SLR.5 — Homoskedasticity.] $\Var{u \given x} = \sigma^2$ (the error variance does not depend on $x$).
\end{description}
}

\thm{Unbiasedness of OLS}{
Under SLR.1--SLR.4, the OLS estimators are unbiased:
\[
\E{\hb_1} = \beta_1, \qquad \E{\hb_0} = \beta_0.
\]
}

\pf{
Write the slope as a function of the errors. Substituting $y_i = \beta_0+\beta_1 x_i + u_i$ into the numerator and using $\sumi(x_i-\bar x)=0$,
\[
\ba
\hb_1 &= \frac{\sumi (x_i-\bar x)\,y_i}{\text{SST}_x}
       = \frac{\sumi (x_i-\bar x)(\beta_0+\beta_1 x_i + u_i)}{\text{SST}_x}\\
      &= \frac{\beta_1 \sumi (x_i-\bar x)(x_i-\bar x) + \sumi (x_i-\bar x)u_i}{\text{SST}_x}
       = \beta_1 + \frac{\sumi (x_i-\bar x)u_i}{\text{SST}_x}.
\ea
\]
Conditioning on the sample of $x$'s and applying SLR.4 ($\E{u_i\given x}=0$),
\[
\E{\hb_1} = \beta_1 + \frac{1}{\text{SST}_x}\sumi (x_i-\bar x)\,\E{u_i \given x} = \beta_1.
\]
For the intercept, $\hb_0 = \bar y - \hb_1\bar x = \beta_0 + (\beta_1-\hb_1)\bar x + \tfrac1n\sumi u_i$, so $\E{\hb_0} = \beta_0$ by the result just shown and $\E{u_i}=0$.
}

If any of SLR.1--SLR.4 fails, unbiasedness generally fails too. The fragile one is SLR.4: whenever an omitted factor inside $u$ is correlated with $x$, the slope is biased --- the central theme of Chapter~\ref{ch:mlr}.

\thm{Sampling Variance of OLS}{
Under SLR.1--SLR.5,
\[
\Var{\hb_1} = \frac{\sigma^2}{\text{SST}_x},
\qquad
\Var{\hb_0} = \frac{\sigma^2\, n^{-1}\sumi x_i^2}{\text{SST}_x}.
\]
}

\pf{
From $\hb_1 = \beta_1 + \dfrac{\sumi (x_i-\bar x)u_i}{\text{SST}_x}$, treating the $x$'s as fixed and using independence across $i$ together with homoskedasticity $\Var{u_i\given x}=\sigma^2$,
\[
\Var{\hb_1} = \frac{1}{(\text{SST}_x)^2}\sumi (x_i-\bar x)^2\,\Var{u_i\given x}
            = \frac{\sigma^2}{\text{SST}_x}.
\]
}

Two messages are already visible in $\Var{\hb_1}=\sigma^2/\text{SST}_x$: the slope is estimated more precisely when the error is small ($\sigma^2$ low) and when the regressor varies a lot ($\text{SST}_x$ large). More data raises $\text{SST}_x$ and so shrinks the variance.

The error variance $\sigma^2$ is unknown and must itself be estimated.

\defn{Error Variance Estimator and Standard Error}{
The unbiased estimator of $\sigma^2$ and the implied standard error of the slope are
\[
\widehat\sigma^2 = \frac{1}{n-2}\sumi \hat u_i^{\,2},
\qquad
\widehat\sigma = \sqrt{\widehat\sigma^2},
\qquad
\operatorname{se}(\hb_1) = \frac{\widehat\sigma}{\sqrt{\text{SST}_x}}.
\]
The divisor $n-2$ corrects for the two parameters ($\beta_0,\beta_1$) already used to form the residuals; $\widehat\sigma$ is called the \emph{standard error of the regression} (SER).
}

\rmkb[Where this is heading]{We have, for one regressor, found the OLS estimator, shown it is unbiased under SLR.1--SLR.4, and derived its variance under the extra homoskedasticity assumption SLR.5. Chapter~\ref{ch:mlr} repeats this exact program with many regressors, where the new and important phenomenon is omitted-variable bias; Chapter~\ref{ch:hetero} returns to drop SLR.5 and asks what survives under heteroskedasticity.}
